% !TeX root = ../main.tex

\section{Assignment \#2}

\begin{problem}[note = Isotope shifts in the rotational spectrum of OCS (2pts),
                label = hw2:problem1]
  \leavevmode \vspace*{-.5\baselineskip}
  \begin{snugshade}
    \begin{enumext}[label = \ensuremath \ast]
      \item
      Consider linear carbonyl sulfide, \ce{^{16}O ^{12}C ^{32}S} with the
      molecular structure \ce{O = C = S}.
      Assume the molecule behaves as a rigid linear rotor.
      \item
      Use the following approximate bond lengths:
      $r_{\ce{C-O}} = \qty{1.16}\angstrom$.
      Set the oxygen atom at $x = 0$, so that $x_{\ce O} = 0$,
      $x_{\ce C} = \qty{1.16}\angstrom$, $x_{\ce S} = \qty{2.72}\angstrom$.
      Use approximate isotrope masses for convenience:
      \ce{^{16}O}: $16$, \ce{^{18}O}: $18$, \ce{^{12}C}: $12$,
      \ce{^{13}C}: $13$, \ce{^{32}S}: $32$, \ce{^{34}S}: $34$.
      \item
      For a linear rigid rotor, $\tilde\nu_{J \to J+1} = 2B(J + 1)$, and for the
      $J = 3 \to 4$ transition, $\tilde \nu_{3 \to 4} = 8B$,
      where $B = \frac{h}{8\pi^2cI}$
      (for \ce{^{16}O ^{12}C ^{32}S}, $B = \qty{0.20307}{\per\cm}$).
    \end{enumext}
  \end{snugshade}
  \paragraph{Questions.}
  \begin{enumext}
    \item If one atom is replaced by an isotope, calculate the shift of the
    $J = 3 \to 4$ line for each isotopologue:
    \[
      \ce{^{18}O ^{12}C ^{32}S}, \quad
      \ce{^{16}O ^{13}C ^{32}S}, \quad
      \ce{^{16}O ^{12}C ^{34}S}.
    \]
    Report the shifts in $\unit{\per\cm}$ relative to the original species.
    \item If the spectrometor cannot resolve frequency shifts smaller than
    $\qty{0.05}{\per\cm}$, which isotope-labeled species cannot be resolved
    from this main peak?
  \end{enumext}
\end{problem}
\begin{solution}\leavevmode
  \begin{enumext}
    \item Due to $\tilde \nu_{J \to J+1} \propto B$
    and $B = \frac{h}{8\pi^2cI}$, we can obtain
    \[
      \frac{\tilde \nu'}{\tilde \nu} = \frac{B'}{B} = \frac{I}{I'}, \quad
      \Delta \tilde \nu = \tilde \nu' - \tilde \nu
    = \tilde \nu \ab(\frac{I}{I'} - 1).
    \]
    In the shift of $J = 3 \to 4$,
    $\Delta \tilde \nu = 8B \ab(\frac{I}{I'} - 1)$.
    The center of mass and the moment of inertia in the original and
    reference cases are
    \[\begin{array}{l@{~}c@{~}l@{\quad}l@{~}c@{~}l@{~}c@{~}l}
      x_c^0     & = & \qty{1.68}\angstrom, &
      I_0       & = & 16 \times (x_c^0)^2 + 12 \times (x_c^0 - 1.16)^2
                    + 32 \times (x_c^0 - 2.72)^2
                & = & \qty{83.0140}{\amu \angstrom\squared},\\
      x_c^i     & = & \qty{1.63}\angstrom, &
      I_1       & = & 18 \times (x_c^0)^2 + 12 \times (x_c^i - 1.16)^2
                    + 32 \times (x_c^i - 2.72)^2
                & = & \qty{88.4940}{\amu \angstrom\squared},\\
      x_c^{ii}  & = & \qty{1.67}\angstrom, &
      I_2       & = & 16 \times (x_c^0)^2 + 13 \times (x_c^{ii} - 1.16)^2
                    + 32 \times (x_c^{ii} - 2.72)^2
                & = & \qty{83.2827}{\amu \angstrom\squared},\\
      x_c^{iii} & = & \qty{1.72}\angstrom, & 
      I_3       & = & 16 \times (x_c^0)^2 + 12 \times (x_c^{iii} - 1.16)^2
                                + 34 \times (x_c^{iii} - 2.72)^2
                & = & \qty{85.0967}{\amu \angstrom\squared}.
    \end{array}\]
    Substituting the data, then the shifts are
    \[
      \Delta \tilde \nu_i     = \qty{-0.1006}{\per\cm}, \quad
      \Delta \tilde \nu_{ii}  = \qty{-0.0052}{\per\cm}, \quad
      \Delta \tilde \nu_{iii} = \qty{-0.0398}{\per\cm}.
    \]
    \item From the results,
    \ce{^{18}O ^{12}C ^{32}S} can be resolved, while \ce{^{16}O ^{13}C ^{32}S}
    and \ce{^{16}O ^{12}C ^{34}S} cannot.
  \end{enumext}
\end{solution}

\pagebreak

\begin{problem}[Strongest rotational line and temperature dependence (2pts)]
  \leavevmode \vspace*{-.5\baselineskip}
  \begin{snugshade}
    \begin{enumext}[label = \ensuremath \ast]
      \item
      For a linear rigid rotor, the approximate intensity of the absorption
      transition $J \to J + 1$ is
      $I_j \propto (J + 1) \exp\ab[-\frac{BhcJ(J + 1)}{k_BT}]$, where $J$ is the
      lower-state rotational quantum number.
      \item
      The strongest transition is approximately determined by
      \[
        J_{\max} \approx \frac{1 + \sqrt{1 + 8k_BT/(Bhc)}}{4} - 1.
      \]
      Use $\frac{k_BT}{hc} = \qty{208.5}{\per\cm}$ at $\qty{300}\K$ and
      $B_{\ce{CO}} = \qty{1.8225}{\per\cm}$.
      \item
      For \ce{OCS}, use the rotational contant obtained
      in~\zcref{hw2:problem1}: $B_{\ce{OCS}} \approx \qty{0.2031}{\per\cm}$.
    \end{enumext}
  \end{snugshade}
  \paragraph{Questions.}
  \begin{enumext}
    \item Calculate the approximate $J_{\max}$ for \ce{CO} at $\qty{300}\K$.
    Which transition is strongest?
    \item Calculate the approximate $J_{\max}$ for \ce{CO} at $\qty{230}\K$.
    Which transition is strongest?
    \item For \ce{OCS} at $\qty{300}\K$, which term in the $J_{\max}$ equation
    changes compared with \ce{CO}? Calculate the approximate $J_{\max}$
    for \ce{OCS}. Which transition is strongest?
  \end{enumext}
\end{problem}
\begin{solution}\leavevmode
  \begin{enumext}
    \item (b)\enspace Substituting the variable into the transition expression
    \[
      J_{\max}(\qty{300}\K)^{\ce{CO}}
    = \frac{1 + \sqrt{1 + 8 \times
            \qty{208.5}{\per\cm}/\qty{1.8225}{\per\cm}}}{4} - 1 = 6.82, \quad
      J_{\max}(\qty{230}\K)^{\ce{CO}} = 5.88.
    \]
    Rounding to the nearest integer, for $\qty{300}\K$, $J_{\max} = 7$,
    The strongest transition is the $J = 7 \to 8$ line;
    For $\qty{230}\K$, $J_{\max} = 6$.
    The strongest transition is the $J = 6 \to 7$ line.
    \stepcounter{enumXi}
    \item The rotational constant $B$ is the term that changes significantly
    compared to \ce{CO}, since \ce{OCS} is heavier than \ce{CO} and they have
    different moment of inertia. The transition is
    \[
      J_{\max}(\qty{300}\K)^{\ce{OCS}}
    = \frac{1 + \sqrt{1 + 8 \times
            \qty{208.5}{\per\cm}/\qty{0.2031}{\per\cm}}}{4} - 1 = 21.91.
    \]
    Rounding to the nearest integer, $J_{\max} = 22$.
    The strongest transition is the $J = 22 \to 23$ line.
  \end{enumext}
\end{solution}

\begin{problem}[Centrifugal distortion for molecular rotation (1pt)]
  \leavevmode \vspace*{-.5\baselineskip}
  \begin{snugshade}
    If the centrifugal distortion is considered, then there is an additional
    contribution of $[-DJ^2 (J + 1)^2]$ to the energetic states.
    Whats is the expression for neighboring rotational line spacing
    $(J \to J + 1)$ under this new condition
    (Ideeal rigid-rotor engines $E_\text{rotation} = hcBJ(J + 1)$)?
  \end{snugshade}
\end{problem}
\begin{solution}
  The energy levels will be written as
  \[
    E(J) = BJ(J + 1) - DJ^2(J + 1)^2.
  \]
  Then, the transition becomes
  \begin{multline*}
    \nu_{J\to J+1} = E(J + 1) - E(J)
    = [B(J + 1)(J + 2) - D(J + 1)^2(J + 2)^2] \\ - [BJ(J + 1) - DJ^2(J + 1)^2]
    = 2B(J + 1) - 4D(J + 1)^3.
  \end{multline*}
\end{solution}